\documentclass[12pt]{article}

\usepackage[T1]{fontenc}
\usepackage{lmodern}
\usepackage{graphicx}
\usepackage{booktabs}
\usepackage{geometry}
\usepackage{setspace}
\usepackage[colorlinks=true, linkcolor=blue, urlcolor=blue, citecolor=blue]{hyperref}
\usepackage{tocloft}
\usepackage{xcolor}
\usepackage{listings}
\usepackage{amsmath}
\usepackage{parskip}
\usepackage{indentfirst}
\usepackage{float}

\geometry{margin=1in}
\onehalfspacing

% Indentation and Paragraph Spacing
\setlength{\parindent}{2em}
\setlength{\parskip}{1em}

\renewcommand{\cftsecleader}{\cftdotfill{\cftdotsep}}

\begin{document}
\pagenumbering{gobble}

% -------- Title Page --------
\begin{titlepage}
    \vspace*{\fill}
    \begin{center}
        \begin{LARGE}
            \textbf{Image to SPICE Netlist Generator using Deep Learning}
        \end{LARGE}
        
        \vspace{10pt}
        
        \begin{Large}
            \textbf{Endterm Report}
        \end{Large}
        
        \vspace{10pt}
        
        \textbf{\Large Course Code : EC304}\\
        \vspace{6pt}
        \textbf{\Large Course Name : Project Phase - I}\\
        
        \vspace{40pt}
        
        \textbf{\large Team Members}\\
        \vspace{10pt}
        \textbf{\large
            Yash Badwaik (U23EC078)\\
            Darshan Nayak (U23EC105)\\
            Ridham Kevat (U23EC121)
        }
        
        \vspace{30pt}
        
        \includegraphics[width=0.3\textwidth]{./logo.png} 
        
        \vspace{25pt}
        
        \textbf{\large
            Guide: Dr. Nithin Chatterji \\
            Co-Guide: Dr. Kishor Upla
        }
        
        \vspace{40pt}
        
        \textbf{\large
            Department of Electronics and Communication Engineering\\
            Sardar Vallabhbhai National Institute of Technology, Surat\\
            April, 2026
        }
    \end{center}
    \vspace*{\fill}
\end{titlepage}

% -------- Front Pages --------
\newpage
\pagenumbering{roman}
\section*{Abstract}
\addcontentsline{toc}{section}{Abstract}
The transition from conceptual hardware sketching to software simulation represents a significant bottleneck in electronic design automation (EDA), traditionally demanding the tedious and manual drafting of SPICE netlists. This project presents an end-to-end, desktop-based machine learning pipeline engineered to automatically digitize hand-drawn circuit schematics into executable NGSPICE simulations. The core of the system leverages a YOLOv8 object detection architecture, fine-tuned via transfer learning to classify nine distinct electronic components while seamlessly handling extreme scale variance. To complement the deep learning model, a deterministic computer vision subsystem utilizes OpenCV for adaptive thresholding and morphological operations, algorithmically routing hand-drawn wire traces onto a discrete coordinate plane. These subsystems are integrated into a custom-built Graphical User Interface (GUI) developed in Python. This interface features an interactive digital canvas, allowing users to dynamically modify AI-detected components prior to the automated generation of the final \texttt{.cir} netlist. Quantitative evaluations, including precision-recall and confusion matrix analyses, demonstrate the model's high classification accuracy and minimal false-positive rate. Ultimately, this system effectively bridges the gap between rapid physical prototyping and rigorous digital simulation, substantially reducing schematic drafting overhead.

\newpage
\section*{Problem Statement}
\addcontentsline{toc}{section}{Problem Statement}
In academic research, rapid electronic prototyping, and analog circuit design, the conversion of visual circuit representations into textual SPICE netlists is a mandatory prerequisite for software simulation. Traditionally, this digitization process relies exclusively on manual data entry within specialized schematic editors. This manual procedure is inherently time-consuming, and the repetitive nature of transcription frequently introduces topological and parameter-value errors that severely hinder the debugging pipeline. Despite significant advancements in artificial intelligence, a robust and accessible interface capable of reliably translating hand-drawn sketches directly into industry-standard SPICE syntax remains largely underdeveloped. 

This project addresses this operational bottleneck by engineering a comprehensive deep learning pipeline. The primary objective is to develop a highly robust system capable of automatically parsing uploaded circuit schematics, accurately localizing and classifying native electronic components, extracting intricate wiring topologies, and synthesizing structurally accurate, NGSPICE-compatible netlists ready for immediate downstream electrical simulation.

\newpage
\section*{Acknowledgements}
\addcontentsline{toc}{section}{Acknowledgements}
We wish to express our profound gratitude to our project guide, \textbf{Dr. Nithin Chatterji}, and co-guide, \textbf{Dr. Kishor Upla}, for their invaluable mentorship, continuous encouragement, and insightful feedback throughout the duration of this project. Their extensive domain expertise and dedicated support have been instrumental in shaping our research direction and resolving the complexities inherent in deep learning model integration and simulation deployment.

We also extend our sincere appreciation to the Department of Electronics and Communication Engineering at the \textbf{Sardar Vallabhbhai National Institute of Technology (SVNIT), Surat}, for providing the rigorous academic environment, computational infrastructure, and essential resources required to execute this research efficiently.

Finally, we extend our thanks to our peers and all individuals who, directly or indirectly, enriched our learning experience and contributed to the successful completion of this engineering project.

% -------- Contents Page --------
\newpage
\tableofcontents

\newpage
\pagenumbering{arabic}

% -------- 1. Introduction --------
\section{Introduction}

The design and simulation of electronic circuits is a fundamental process in electrical engineering, serving as the bridge between theoretical concepts and physical hardware. Traditionally, the ideation phase begins with rapid, hand-drawn sketches on paper or whiteboards. However, to validate these designs mathematically, engineers must manually translate these sketches into simulation software using specific syntax, such as SPICE (Simulation Program with Integrated Circuit Emphasis) netlists. This manual transcription is notoriously tedious, time-consuming, and highly susceptible to human error. As artificial intelligence and computer vision technologies advance, there is a significant opportunity to automate this transition. This project introduces a comprehensive machine learning pipeline designed to seamlessly digitize hand-drawn electronic schematics and convert them directly into executable NGSPICE simulations.

\subsection{Motivation and Objectives}
\noindent \textbf{Motivation}\par
The primary motivation for this project stems from the friction experienced during the hardware prototyping phase. While modern Electronic Design Automation (EDA) tools are powerful, they impose a rigid workflow that interrupts the natural, fluid process of hand-drawing a circuit. For students, educators, and rapid-prototyping engineers, the requirement to manually place components and route wires in a digital GUI—or worse, write raw \texttt{.cir} text files—acts as a severe bottleneck. By applying Deep Learning to visual recognition, we aim to eliminate this bottleneck, allowing engineers to retain the speed of physical sketching while immediately gaining the analytical power of digital simulation.

\vspace{0.5cm}
\noindent \textbf{Objectives}\par
To achieve this, the project was guided by the following core objectives:
\begin{enumerate}
    \item \textbf{Object Detection:} To train and deploy a Convolutional Neural Network (YOLOv8) capable of recognizing and localizing hand-drawn electronic symbols with high precision, overcoming variations in sketching styles and extreme scale differences.
    \item \textbf{Computer Vision Routing:} To implement deterministic image processing algorithms (using OpenCV) that can extract, repair, and trace hand-drawn wire connections between detected components.
    \item \textbf{Interactive GUI Development:} To engineer a custom desktop environment (``PySpice Studio'') that provides users with an interactive digital canvas to visualize the AI's predictions and manually adjust the circuit before simulation.
    \item \textbf{Algorithmic Netlist Generation:} To develop a mathematical routing algorithm that translates the spatial coordinates of the digital canvas into a syntactically correct NGSPICE netlist, executing the simulation and retrieving the results seamlessly.
\end{enumerate}

\subsection{Scope of the project (Desktop-based ML pipeline)}

It is crucial to define the boundaries of this project to accurately evaluate its success. This system is designed as an end-to-end, \textbf{desktop-based machine learning pipeline} built entirely in Python.

\vspace{0.5cm}
\noindent \textbf{In-Scope Features:}
\begin{itemize}
    \item The system is capable of recognizing \textbf{9 distinct classes} of electronic components (including standard passives, basic active components like BJTs, and power sources).
    \item The pipeline operates locally on the user's machine, integrating the YOLOv8 inference engine, the OpenCV processing logic, and the NGSPICE binary execution within a unified Tkinter graphical interface.
    \item The system supports the algorithmic extraction of orthogonal (Manhattan) wire routing to generate node-accurate netlists.
\end{itemize}

% -------- 2. Research Overview Done by Us --------
\section{Research and Literature Review}

Before developing the custom PySpice Studio pipeline, an extensive review of existing tools, datasets, and machine learning models was conducted. This phase was critical for understanding the limitations of current electronic design automation (EDA) workflows and selecting the optimal technology stack for real-time circuit digitization.

\subsection{Finding the dataset (Kaggle circuit dataset and its limitations)}

Training a robust Deep Learning model requires a highly diverse and accurately annotated dataset. For this project, we sourced a specialized dataset from Kaggle containing images of hand-drawn electronic circuits (see Figure~\ref{fig:dataset_samples}). The dataset encompasses 19 distinct classes, ranging from basic passive components (resistors, capacitors, inductors) to active components (MOSFETs, BJTs, Op-Amps) and structural elements (junctions, wires, and text).

While the dataset provided a strong foundational baseline, our research identified several limitations that had to be addressed during the preprocessing phase:
\begin{itemize}
    \item \textbf{Class Imbalance:} Common components like resistors and wires were heavily overrepresented compared to complex components like transformers or logic gates.
    \item \textbf{Illumination Variance:} The images contained significant variations in lighting, shadows, and paper textures, mimicking real-world conditions but complicating feature extraction.
    \item \textbf{Lack of Value Annotation:} While the dataset localized text (e.g., drawing a bounding box around "10k"), it did not transcribe the actual strings. This limitation scoped our current project to component topology extraction, delegating Optical Character Recognition (OCR) to future work.
\end{itemize}


\begin{figure}[H]  % figure is at wrong place in pdf it should be in finding dataset
    \centering
     \includegraphics[width=0.7\textwidth]{sample_from_dataset.jpg}
    \caption{Sample images from the Kaggle dataset showcasing varying sketching styles, illumination, and component orientations.}
    \label{fig:dataset_samples}
\end{figure}

\subsection{Review of existing object detection models}

Object detection is a mature field within computer vision, but not all models are suited for schematic extraction. We reviewed several standard architectures:
\begin{itemize}
    \item \textbf{Template Matching (OpenCV):} Traditional computer vision relies on sliding a template image across the input. Research showed this method is highly brittle; it fails immediately if the user draws a resistor slightly larger, rotated, or messier than the template.
    \item \textbf{Faster R-CNN:} A two-stage Deep Learning detector. It first generates "region proposals" and then classifies them. While highly accurate, the two-stage pipeline is computationally heavy and slower, making it less ideal for a responsive desktop application.
    \item \textbf{Single Shot MultiBox Detector (SSD):} Faster than R-CNN as it predicts bounding boxes in a single pass, but historically struggles with detecting very small objects (like junction dots in a large circuit diagram).
\end{itemize}

\subsection{Exploring schematic drawing tools (SchemDraw)}

To understand how software programmatically generates circuit diagrams, we researched existing Python libraries, primarily \texttt{SchemDraw}. \texttt{SchemDraw} allows engineers to draw highly professional schematics by writing Python code (e.g., \texttt{d += elm.Resistor().right().label('10k')}). 

Our research into \texttt{SchemDraw} was invaluable for understanding how electronic symbols are rendered on a digital coordinate plane (an example output is provided in Figure~\ref{fig:schemdraw}). However, it confirmed the core problem statement: writing lines of code to generate a visual diagram is inherently backward for a hardware prototyping workflow. The engineer must already know the exact topology and calculate the grid coordinates mentally. This solidified our goal to build a visual-first, drag-and-drop canvas (PySpice Studio) rather than a code-first renderer.

% �� IMAGE SUGGESTION: Put a side-by-side image here. Left side: 3 lines of SchemDraw python code. Right side: The circuit it draws.
\begin{figure}[H]
    \centering
     \includegraphics[width=0.8\textwidth]{images/schemdraw_example.png}
    \caption{An example of programmatic circuit generation using the SchemDraw library, demonstrating the syntax-heavy nature of manual digitization.}
    \label{fig:schemdraw}
\end{figure}

\subsection{Understanding the NGSPICE simulation standard}

To bridge the gap between computer vision and physical simulation, an in-depth study of the SPICE (Simulation Program with Integrated Circuit Emphasis) standard was required. SPICE does not understand graphics, lines, or bounding boxes; it only understands nodal matrices represented as plain text (Netlists).

A standard NGSPICE netlist entry follows a strict syntactical format: \texttt{[Component\_Name] [Node\_1] [Node\_2] [Value]}. 
Our research revealed that the true complexity of this project would not be identifying the components, but rather mapping the spatial interconnects (the hand-drawn wires) into discrete, shared mathematical nodes. For example, if a resistor and a capacitor touch the same wire, the Python backend must assign them an identical string identifier (e.g., \texttt{N1}) before invoking the NGSPICE binary.

\subsection{Finding the perfect image processing model (YOLOv8)}

After evaluating the limitations of traditional OpenCV methods and two-stage detectors, we concluded that the YOLO (You Only Look Once) architecture—specifically YOLOv8—was the optimal solution for this project. 

Unlike two-stage models that first generate region proposals before classification, YOLOv8 consolidates these distinct theoretical steps. It frames object detection as a single regression problem and utilizes an anchor-free detection head. This is highly advantageous for electronic components that vary wildly in aspect ratio (e.g., a long, thin wire versus a perfectly square logic gate). Furthermore, YOLOv8's integration with PyTorch allows for seamless exporting of the trained weights (\texttt{.pt} files) directly into our local Python Tkinter environment, enabling rapid, offline inference without the need for cloud computing APIs.

% �� IMAGE SUGGESTION: A simple block diagram showing the difference between Two-Stage (R-CNN) and One-Stage (YOLO) detectors.


% -------- 3. About ML Model (YOLO \& Computer Vision) --------
\section{Machine Learning Architecture and Vision Pipeline}

To successfully convert a raw, hand-drawn sketch into a functional digital simulation, the system relies on a hybrid pipeline combining cutting-edge Deep Learning object detection with traditional, deterministic Computer Vision algorithms. This chapter details the theoretical and mathematical foundations of the vision subsystem.

\subsection{The evolution of Deep Learning in image processing}

Traditional image processing relies heavily on manual feature extraction techniques, such as Haar cascades or template matching, to identify objects. While effective in highly controlled environments, these methods struggle severely with varying lighting, scale discrepancies, and the inherent inconsistencies of human handwriting. Deep Learning, specifically Convolutional Neural Networks (CNNs), revolutionized this domain by automatically learning hierarchical features from massive datasets.

Early deep learning detection architectures, such as Faster R-CNN, utilized a two-stage process: they first generated thousands of "region proposals" and subsequently ran a classifier over each one. While highly accurate, this two-stage pipeline is computationally expensive. For real-time circuit digitization, our system requires rapid, unified inference. Therefore, we selected the \textbf{YOLO (You Only Look Once)} architecture. YOLO frames object detection as a single regression problem, predicting bounding boxes and class probabilities simultaneously directly from the full image in one evaluation.

\subsection{What is YOLOv8 and how it works (Backbone, Neck, Head)}

YOLOv8 represents the state-of-the-art in the YOLO family, offering an anchor-free detection mechanism. The architecture is fundamentally divided into three core deep learning subsystems: the Backbone, the Neck, and the Head, as illustrated in Figure~\ref{fig:yolo_architecture}.

\subsubsection{The Backbone (Feature Extractor)}
The backbone is responsible for extracting visual features from the raw input image of the circuit sketch. YOLOv8 utilizes a modified CSPNet (Cross Stage Partial Network). As the image passes through successive convolutional layers, spatial dimensions are reduced while the channel depth increases. Early layers detect primitive features like straight lines and corners (crucial for wires), while deeper layers detect complex semantic features, such as the specific zig-zag pattern of a resistor.

\subsubsection{The Neck (Feature Fusion)}
Circuit diagrams contain objects of wildly different scales—for instance, a massive voltage source versus a microscopic junction dot. To solve this, YOLOv8 uses a PANet (Path Aggregation Network) structure in its Neck. It takes high-level, semantically strong features from the deep layers and combines them with low-level, spatially accurate features from earlier layers. This multi-scale fusion ensures that the model does not lose the coordinates of tiny components while processing the larger image context.

\subsubsection{The Head (Anchor-Free Detection)}
Unlike older models that used predefined "anchor boxes," YOLOv8 directly predicts the center of an object and the distance to its bounding box edges. The head outputs two separate streams simultaneously:
\begin{itemize}
    \item \textbf{Classification Branch:} Predicts the probability of the object belonging to one of the target classes.
    \item \textbf{Regression Branch:} Predicts the spatial bounding box coordinates $[x, y, w, h]$.
\end{itemize}

\begin{figure}[H]
    \centering
    % \includegraphics[width=0.8\textwidth]{images/yolo_architecture.png}
    \caption{High-level architectural block diagram of YOLOv8, illustrating the flow of tensors through the Backbone, PANet Neck, and Dual-Branch Head.}
    \label{fig:yolo_architecture}
\end{figure}

\subsection{The mathematics behind YOLO (IoU and Loss functions)}

During training and inference, the neural network relies on strict mathematical formulations to refine its accuracy and penalize incorrect guesses.

\subsubsection{Intersection over Union (IoU)}
To determine if a predicted bounding box correctly captures a circuit component, the model calculates the IoU against the ground truth labels provided in the dataset. It is defined as the area of overlap divided by the area of union between the predicted box and the actual box:
\begin{equation}
    IoU = \frac{\text{Area of Overlap}}{\text{Area of Union}}
\end{equation}
If the IoU exceeds a predefined threshold (e.g., $0.50$), the detection is considered a True Positive.

\subsubsection{The Loss Function}
During the training phase, YOLOv8 optimizes its weights using a composite loss function that penalizes classification errors, bounding box coordinate errors, and distribution focal loss (DFL) for fine-tuning box edges:
\begin{equation}
    \mathcal{L}_{total} = \lambda_{box} \mathcal{L}_{box} + \lambda_{cls} \mathcal{L}_{cls} + \lambda_{dfl} \mathcal{L}_{dfl}
\end{equation}
By minimizing $\mathcal{L}_{total}$ over hundreds of epochs, the neural network converges, learning the precise visual signatures of the 19 distinct electronic classes.

\subsection{OpenCV image preprocessing (Grayscale, Blur, Adaptive Thresholding)}

While YOLO handles the classification of distinct components, extracting the hand-drawn wires requires deterministic pixel-level analysis. Before the system can mathematically route wires, the raw image must be cleaned using OpenCV to remove shadows, paper textures, and lighting gradients (Figure~\ref{fig:preprocessing}).

\begin{itemize}
    \item \textbf{Grayscale Conversion:} The image is reduced from a 3-channel RGB array to a 1-channel grayscale matrix, reducing computational overhead as color is irrelevant to circuit topology.
    \item \textbf{Gaussian Blurring:} A $5 \times 5$ Gaussian kernel is convoluted over the image to smooth out high-frequency noise and paper grain.
    \item \textbf{Adaptive Thresholding:} Because a photograph of a sketch often has uneven lighting, standard global thresholding fails. We utilized Gaussian Adaptive Thresholding, which calculates the threshold for small, localized regions of the image, ensuring that dark lines are extracted cleanly regardless of shadows.
\end{itemize}

\begin{figure}[H]
    \centering
    % \includegraphics[width=1.0\textwidth]{vision_output/1_preprocessing_pipeline.png}
    \caption{The OpenCV preprocessing pipeline converting a raw photograph into a clean binary matrix suitable for wire extraction.}
    \label{fig:preprocessing}
\end{figure}

\subsection{Morphological operations for wire enhancement}

Hand-drawn circuit sketches frequently suffer from discontinuities; a pen might skip, leaving a micro-gap in a wire trace. If left uncorrected, the routing algorithm would perceive this as an open circuit, ruining the final SPICE simulation.

To repair these breaks algorithmically, we employed Morphological Operations on the binarized image using a $5 \times 5$ structural element (kernel).
\begin{itemize}
    \item \textbf{Dilation:} This operation expands the white regions (the wire pixels) of the image, effectively bridging small gaps, but thickening the overall lines.
    \item \textbf{Morphological Closing:} To bridge gaps without permanently thickening the lines, we used the "Closing" operation, which is a mathematical Dilation immediately followed by an Erosion. This successfully repairs broken traces while preserving the original structural geometry of the circuit.
\end{itemize}

By actively combining these filtering techniques, the system mathematically ensures that any pen skips and micro-gaps drawn by the user are repaired and do not register as false open circuits during the final netlist generation phase.

% -------- 4. How YOLO is Used in Our Project --------
\section{How YOLO is Used in Our Project}

While understanding the base YOLO architecture is critical, the true engineering challenge of this project lay in adapting a generalized Deep Learning model into a specialized electronic schematic digitizer. This section details the practical implementation, training, and data-translation phases of the machine learning pipeline.

\subsection{Domain-Specific Fine-Tuning for 19 Electronic Components}

A base YOLOv8 model, pre-trained on the standard COCO (Common Objects in Context) dataset, is highly proficient at recognizing everyday objects like cars, people, and furniture, but it has zero concept of electronic symbology. Training a deep neural network from scratch to recognize circuits would require millions of images and massive computational resources.

To circumvent this, we utilized \textbf{Transfer Learning}. We retained the generalized feature-extraction weights in the YOLOv8 backbone—which already understood how to detect edges, corners, and geometric shapes—but completely replaced the classification head. The model was then fine-tuned on our specialized Kaggle dataset containing 19 unique circuit classes (including resistors, capacitors, voltage sources, ground symbols, etc.). This approach allowed the network to leverage its existing foundational knowledge of geometry while exclusively learning the visual signatures of electronic components.

\subsection{Handling Extreme Scale Variance in Schematics}

One of the primary difficulties in parsing hand-drawn schematics is the extreme variance in object scale. A long wire trace might span $80\%$ of the image canvas, while a junction dot or a small diode takes up less than $2\%$. Standard convolutional networks often lose the spatial data of tiny objects as the image is pooled into deeper layers.

We explicitly leveraged YOLOv8’s multi-scale PANet (Path Aggregation Network) neck to solve this. By extracting bounding box predictions from feature maps at three different resolution scales (typically $80 \times 80$, $40 \times 40$, and $20 \times 20$), the model maintains situational awareness. The high-resolution feature maps successfully localized the microscopic components (like text and junctions), while the low-resolution, high-semantic maps accurately bounded massive components (like transformers and long wire segments).

\subsection{Model training parameters and Google Colab setup}

Due to the heavy computational requirements of training a deep neural network, the model was trained using cloud infrastructure via Google Colab, utilizing a hardware-accelerated GPU environment.

The training process was executed using the following hyperparameter configuration:
\begin{itemize}
    \item \textbf{Image Size (\texttt{imgsz}):} 640x640 pixels. All input images were dynamically scaled and padded to this dimension to maintain a uniform tensor shape.
    \item \textbf{Epochs:} The model was trained for 100 iterations over the entire dataset. Early stopping was monitored to prevent overfitting.
    \item \textbf{Batch Size:} 16 images per forward/backward pass, optimizing GPU memory utilization.
    \item \textbf{Optimizer:} SGD (Stochastic Gradient Descent) with momentum, allowing the network to navigate the loss landscape efficiently.
\end{itemize}
Upon completion of the training phase, the optimized weights were exported as a \texttt{best.pt} file. This highly compressed, standalone file contains the entire "brain" of our AI, allowing it to be downloaded and run locally on the PySpice Studio desktop environment without requiring an internet connection. Figure~\ref{fig:training_results} displays the training metrics over 100 epochs.

% �� IMAGE SUGGESTION: Insert the "results.png" generated by Colab here. It's a grid of graphs showing Loss going down and mAP going up!
\begin{figure}[H]
    \centering
    % \includegraphics[width=1.0\textwidth]{runs/detect/train/results.png}
    \caption{Training metrics over 100 epochs, demonstrating the steady decrease in bounding box and classification loss, alongside the corresponding rise in Mean Average Precision (mAP).}
    \label{fig:training_results}
\end{figure}

\subsection{Bridging the gap: Translating AI tensors to pixel coordinates}

The most critical integration point of the project occurs when the Deep Learning pipeline hands its data over to the deterministic Logic Pipeline (PySpice Studio). When a user uploads a hand-drawn sketch to the desktop interface, the image is passed through the \texttt{.predict()} method of our trained \texttt{best.pt} model.

The direct output of YOLO is not a visual image, but rather a multi-dimensional mathematical tensor. For every detected object, the tensor contains an array formatted as:
\begin{equation}
    [\text{Class\_ID}, \text{Confidence\_Score}, x_{center}, y_{center}, \text{width}, \text{height}]
\end{equation}

Our Python implementation acts as the bridge. It parses this raw tensor array with the following logic:
\begin{enumerate}
    \item \textbf{Thresholding:} It filters out any detection with a \texttt{Confidence\_Score} below $0.40$ (40\%) to prevent visual hallucinations and false positives.
    \item \textbf{Coordinate Translation:} It translates the normalized center coordinates $(x_c, y_c)$ into absolute pixel coordinates representing the top-left $(x_1, y_1)$ and bottom-right $(x_2, y_2)$ corners of the bounding box.
    \item \textbf{Data Injection:} These coordinates, mapped to their specific string labels (e.g., "Resistor"), are injected into the PySpice Studio state manager.
\end{enumerate}

This translation converts static pixels into interactive data objects—allowing the custom GUI to draw selectable, movable components precisely over the user's original sketch (Figure~\ref{fig:tensor_translation}).

% �� IMAGE SUGGESTION: Show a raw sketch next to the AI's bounding box output.
\begin{figure}[H]
    \centering
    % \includegraphics[width=0.9\textwidth]{images/tensor_translation.png}
    \caption{The visual result of translating YOLO coordinate tensors back into the 2D image plane, highlighting components with labeled bounding boxes.}
    \label{fig:tensor_translation}
\end{figure}
% -------- 5. About Implementation (PySpice Studio) --------
\section{System Implementation: PySpice Studio}

While the deep learning pipeline is responsible for recognizing the electronic components, the system requires a robust logic engine to visualize the results, allow user modifications, and translate visual data into mathematical models. This chapter details the implementation of \textbf{PySpice Studio}, a custom-built, desktop-based Electronic Design Automation (EDA) environment developed entirely in Python.

\subsection{Graphical user interface (GUI) design using Tkinter}

The user interface of PySpice Studio was engineered using Tkinter, Python's standard GUI framework. Tkinter was selected for its lightweight footprint and native OS integration, which ensures fast rendering without the overhead of heavy web frameworks like Electron. The interface layout is presented in Figure~\ref{fig:pyspice_editor}.

The application architecture is divided into three primary regions:
\begin{enumerate}
    \item \textbf{Dynamic Toolbar:} Located at the top, it parses the component database (\texttt{library.py}) to dynamically generate tool buttons organized by categories (Passives, Sources, Active components, and Drawing Tools). 
    \item \textbf{Interactive Canvas:} The central viewport where the circuit schematics are rendered and manipulated.
    \item \textbf{Property Inspector Sidebar:} A dynamic data-binding panel on the right side. When a user selects a component, the sidebar dynamically generates entry fields for the component's internal dictionary (\texttt{comp.params}), allowing users to easily modify reference designators (e.g., changing "R1" to "R2") and SPICE parameters (e.g., changing "1k" to "10k").
\end{enumerate}

\begin{figure}[H]
    \centering
    \includegraphics[width=1.0\textwidth]{images/editor.png}
    \caption{The PySpice Studio graphical user interface displaying a mixed-signal circuit layout, component properties sidebar, and the dynamic toolbar.}
    \label{fig:pyspice_editor}
\end{figure}

\subsection{Interactive canvas engine (panning, zooming, and drag-and-drop)}

To provide a professional CAD-like experience, the PySpice Studio canvas engine implements custom affine transformation logic rather than relying on static Tkinter pixel coordinates. The system maintains two parallel coordinate systems: \textbf{World Coordinates} (the absolute mathematical position of the components) and \textbf{Screen Coordinates} (where the pixels are drawn on the monitor).

All mouse events are piped through \texttt{world\_to\_screen} and \texttt{screen\_to\_world} transformation functions:
\begin{itemize}
    \item \textbf{Grid Snapping:} As the user drags the mouse, coordinates are divided by a constant \texttt{GRID\_SIZE} (set to 20), rounded, and multiplied back. This strictly enforces orthogonal alignment, which is critical for successful netlist generation.
    \item \textbf{Zooming:} Bound to the mouse wheel, the zoom function applies a scalar multiplier (\texttt{self.zoom}) to the canvas and dynamically adjusts the viewing offset (\texttt{offset\_x}, \texttt{offset\_y}) to ensure the zoom centers precisely on the user's cursor.
    \item \textbf{Image Caching:} To maintain a high frame rate during panning, component graphics (rendered via the \texttt{schemdraw} library) are cached in memory using their type, rotation angle, and current zoom scale as a unique dictionary key.
\end{itemize}

\subsection{The wire routing algorithm (converting OpenCV masks to mathematical lines)}

A circuit is fundamentally a graph, and the wires represent the edges connecting the component nodes. Within PySpice Studio, wires are not stored as graphical pixels; they are stored as mathematical tuples containing start and end coordinates: \texttt{((x1, y1), (x2, y2))}.

When the Deep Learning and OpenCV pipelines process a hand-drawn sketch, the morphological line-tracing algorithm (discussed in Section 3) outputs a series of coordinate paths. These paths are passed into the canvas engine, where they are discretized and snapped to the nearest \texttt{GRID\_SIZE} intervals. 

Furthermore, the GUI allows for manual wire correction. Using orthogonal routing logic, when a user clicks two points on the canvas, the system automatically resolves the diagonal delta into horizontal and vertical segments. If a wire intersects a component's registered pin coordinates (defined in \texttt{PIN\_MAP}), an electrical connection is successfully established in the system's memory.

\subsection{Netlist generation logic (converting canvas coordinates to a .cir file)}

SPICE simulators do not understand 2D coordinates; they only understand shared nodes. Translating the visual canvas into a text-based \texttt{.cir} file requires an algorithmic graph traversal, which is handled by the \texttt{analyze\_circuit} function in the \texttt{netlist.py} module.

The generation process follows a strict sequence:
\begin{enumerate}
    \item \textbf{Adjacency Mapping:} The algorithm iterates over all wire segments and component pins, building an adjacency list (\texttt{adj}) that maps every coordinate point to all other points it physically touches.
    \item \textbf{Node Clustering (DFS):} Using a Depth-First Search (DFS) algorithm with a \texttt{visited} set and a \texttt{stack}, the system traverses the adjacency list. Every connected cluster of wires and pins is assigned a unique string ID (e.g., "1", "2"). 
    \item \textbf{Ground and Label Overrides:} If the DFS traversal encounters a Ground symbol (\texttt{gnd}), that entire node cluster is aggressively renamed to the universal SPICE ground identifier: \texttt{"0"}. Custom user labels (e.g., "OUT") overwrite the numeric IDs similarly.
    \item \textbf{Netlist Transcription:} Finally, the algorithm iterates through the active components. It looks up the mathematical nodes touching its pins, queries the component database for its specific SPICE syntax string (e.g., \texttt{\{name\} \{n1\} \{n2\} \{value\}} for a resistor), and formats the final string block.
\end{enumerate}

\subsection{Executing NGSPICE binary simulations directly from Python}

The ultimate goal of the pipeline is mathematical validation. PySpice Studio includes a dedicated Simulation Dialog that allows users to configure Transient, AC, DC Sweep, or Operating Point analyses. 

Once the simulation parameters and plot probes are configured, the Python backend executes the following steps:
\begin{enumerate}
    \item It writes the generated string data from the netlist algorithm into a temporary local file named \texttt{circuit.cir}.
    \item It appends a SPICE \texttt{.control} block to the file, which dictates plot colors, variables to trace, and the specific \texttt{run} command.
    \item It utilizes Python's \texttt{subprocess.Popen} library to asynchronously invoke the \texttt{ngspice.exe} binary installed on the host machine, passing the \texttt{.cir} file as a command-line argument.
\end{enumerate}

This opens the native NGSPICE graphing utility directly over the PySpice Studio interface, providing the user with immediate, highly accurate waveform visualizations of the circuit they originally sketched on paper (Figures~\ref{fig:sim_transient} and \ref{fig:sim_ac}).

\begin{figure}[H]
    \centering
    \includegraphics[width=1\textwidth]{images/transient_rectifier.png}
    \caption{Transient analysis of a Diode Rectifier circuit executed via NGSPICE.}
    \label{fig:sim_transient}
\end{figure}

\begin{figure}[H]
    \centering
    \includegraphics[width=1\textwidth]{images/AC_HPF.png}
    \caption{AC Frequency sweep analysis of a High-Pass Filter plotted on a logarithmic scale.}
    \label{fig:sim_ac}
\end{figure}

% -------- 6. Current Results --------
\section{Current Results}

At the current stage, the project can confidently locate and extract primary components from varying layouts of handwritten circuit designs. We have verified the system output against multiple moderately complex test circuits, transitioning them from physical sketches to mathematical representations (e.g., the raw sketch in Figure~\ref{fig:raw_input}).

% �� PLACEHOLDER: Raw Input Sketch
\begin{figure}[H]
    \centering
    \includegraphics[width=1\textwidth]{C171_D1_P1/C171_D1_P1.jpeg}
    \caption{The raw, hand-drawn input sketch before any digital processing or AI inference.}
    \label{fig:raw_input}
\end{figure}

The feasibility of automated circuit digitization has proven successful. The YOLOv8 model demonstrates robust bounding box captures regardless of the drawing scale or minor rotational variances, as seen in Figure~\ref{fig:example_output}. Figure~\ref{fig:heatmap} provides the corresponding heatmap overlay.

\begin{figure}[H]
    \centering
    \includegraphics[width=1\textwidth]{C171_D1_P1/3_feature_extraction.png}
    \caption{Feature extraction and detection output, demonstrating robust bounding box captures over the raw sketch.}
    \label{fig:example_output}
\end{figure}

\begin{figure}[H]
    \centering
    \includegraphics[width=1\textwidth]{C171_D1_P1/0_yolo_heatmap.png}
    \caption{Heatmap overlay representing the model's detection confidence. Hotter (red) zones indicate near-certainty of a component's center point.}
    \label{fig:heatmap}
\end{figure}

\subsection{Evaluation metrics (Precision, Recall, F1-Score, mAP)}

To scientifically quantify the performance of the object detection model, we rely on four standard machine learning metrics calculated by comparing the AI's predictions against the ground-truth labeled dataset:

\begin{itemize}
    \item \textbf{Precision:} Measures the accuracy of the positive predictions. Out of all the bounding boxes the AI drew, what percentage were actually correct?
    \begin{equation}
        Precision = \frac{True Positives}{True Positives + False Positives}
    \end{equation}
    
    \item \textbf{Recall:} Measures the model's ability to find all the positive samples. Out of all the actual components drawn on the paper, what percentage did the AI successfully find?
    \begin{equation}
        Recall = \frac{True Positives}{True Positives + False Negatives}
    \end{equation}
    
    \item \textbf{F1-Score:} The harmonic mean of Precision and Recall. Since forcing the AI to be highly precise usually lowers its recall (and vice versa), the F1-Score provides a single metric to judge the perfect balance.
    
    \item \textbf{mAP (Mean Average Precision):} The most critical metric for object detection. It is calculated by taking the area under the Precision-Recall curve for every single class, and averaging them together. An mAP@0.5 evaluates the accuracy assuming a bounding box overlap (IoU) of at least 50\% is considered a success.
\end{itemize}

\subsection{Training performance analysis}

The model was trained over 100 epochs, allowing the neural network ample iterations to adjust its internal weights. During the early epochs, the bounding box regression loss and classification loss dropped sharply, indicating that the model quickly learned to differentiate between blank paper, wires, and actual components. 

By epoch 80, the loss curves began to plateau, indicating convergence. The model successfully learned the visual signatures of the 19 classes without overfitting to the training data, as validated by the steady, parallel climb of the validation mAP metrics.

\subsection{Confusion Matrix analysis (Where our model excels and struggles)}

The confusion matrix is a powerful diagnostic tool that visualizes exactly where the AI is making classification errors. In a normalized confusion matrix, the diagonal from top-left to bottom-right represents correct predictions (True Positives).

\begin{figure}[H]
    \centering
    \includegraphics[width=1\textwidth]{images/confusion_matrix_normalized.png}
    \caption{Normalized confusion matrix of the detection model across the 9 component classes and the background.}
    \label{fig:confusion_matrix_normalized}
\end{figure}

\textbf{Analysis of the Matrix:}\par
The deep blue shading along the main diagonal indicates a high overall classification accuracy for most core components, as illustrated in Figure~\ref{fig:confusion_matrix_normalized}. The model performs exceptionally well at identifying geometrically distinct features such as \textbf{Transistors} (88\% accuracy) and \textbf{Resistors} (82\% accuracy). 

However, analyzing the off-diagonal elements reveals specific morphological struggles. \textbf{Inductors} proved to be the most difficult class to detect correctly (only 60\% accuracy), with the model frequently confusing true resistors for inductors (12\% of the time). Similarly, there is a minor confusion where true capacitors are misclassified as sources (9\%).

The most significant insights come from the background classifications (False Positives and False Negatives). The bottom row indicates that the model occasionally misses real components entirely (predicting them as background), particularly variable or small symbols like ground (20\% miss rate) and inductors (19\%). Conversely, the rightmost column reveals a notable False Positive issue with the \textbf{text} class: when the model hallucinates a prediction on an empty background, 73\% of the time it incorrectly assumes it is seeing text. This suggests that paper texture, smudges, or stray pen marks are heavily triggering the text detection features.
\subsection{Precision-Recall curves}

The Precision-Recall (PR) curve plots precision against recall for varying confidence thresholds. A perfect model would have a line that hugs the top-right corner of the graph.

\begin{figure}[H]
    \centering
    \includegraphics[width=0.7\textwidth]{images/BoxPR_curve.png}
    \caption{Precision-Recall curve of the YOLO model, achieving an overall mAP@0.5 of 0.837.}
    \label{fig:pr_curve}
\end{figure}

As shown in Figure \ref{fig:pr_curve}, the overall \textbf{mAP@0.5 reached 0.837}, proving the system is highly viable for real-world usage. Breaking down the classes, the model achieved exceptional Average Precision for Transistors (0.916) and Resistors (0.878). The lowest performing class was the Inductor (0.751), likely due to the highly variable loop count and chaotic sketching styles used by different individuals to represent coils.

Furthermore, analyzing the F1-Confidence curve generated during training revealed an optimal peak of \textbf{0.80 at a confidence threshold of 0.343}. Based on this mathematical peak, our PySpice Studio Python backend was explicitly hardcoded to use a confidence threshold of \texttt{0.40} during inference. This slightly stricter threshold deliberately sacrifices a small amount of recall to prioritize high precision, ensuring the UI is not flooded with "hallucinated" components.

\subsection{System walkthrough (Hand-drawn sketch $\rightarrow$ PySpice Canvas $\rightarrow$ Simulation output)}

To demonstrate the full pipeline, we trace a circuit from paper to simulation.

\textbf{Step 1: Digitization.} The user uploads a photograph of a hand-drawn circuit. The YOLOv8 model runs inference, isolating the components, while the OpenCV morphological algorithms extract the wire traces (Figure~\ref{fig:canvas_confirmation}).

\begin{figure}[H]
    \centering
     \includegraphics[width=0.8\textwidth]{C171_D1_P1/canvas_confirmation.png}
    \caption{Confirmation of the digitized elements before rendering them onto the PySpice Studio interactive canvas.}
    \label{fig:canvas_confirmation}
\end{figure}

\textbf{Step 2: Interactive Canvas.} The PySpice Studio GUI translates these raw matrices into an interactive, grid-snapped digital canvas (Figure~\ref{fig:canvas_output}). 

% �� PLACEHOLDER: Canvas Output
\begin{figure}[H]
    \centering
     \includegraphics[width=1.0\textwidth]{C171_D1_P1/canvas_output.png}
    \caption{The final digitized circuit loaded into the PySpice Studio interactive canvas, ready for user modification and wire-routing validation.}
    \label{fig:canvas_output}
\end{figure}

\textbf{Step 3: SPICE Execution.} Once the user verifies the layout, the Python logic subsystem generates the nodal netlist and invokes the NGSPICE binary. The resulting simulation (e.g., a transient wave or AC frequency sweep) is immediately plotted on the screen, successfully completing the journey from a physical sketch to a validated digital simulation.

% -------- 7. Future Work --------
\section{Future Work}

While the current iteration of PySpice Studio successfully bridges the gap between hand-drawn sketches and functional NGSPICE simulations, there remain several avenues for optimization. Future development will focus on enhancing the topological routing logic and expanding the system's analytical capabilities.

\subsection{Implementation of OCR text extraction}

Currently, the system isolates text bounding boxes but assigns default parametric values (e.g., $1\text{k}\Omega$) that require manual GUI editing. Future iterations will integrate an OCR engine (such as EasyOCR) to transcribe handwritten values (e.g., "10k", "5V") and algorithmically associate them with the nearest electronic component. This spatial mapping will fully automate the SPICE parameterization process.

\subsection{Web and Cloud deployment}

To increase accessibility and eliminate the hardware constraints of running local deep learning models, the monolithic Python/Tkinter architecture will be refactored into a client-server web application. Moving YOLO inference and NGSPICE execution to a cloud-hosted FastAPI backend will allow users to digitize circuits natively from mobile devices and web browsers.

\subsection{Source classification}

The AI model occasionally struggles to differentiate between visually similar circular source symbols if the sketch is blurry. Future work will implement a secondary, high-resolution classification pass to reliably distinguish between highly specific inputs, such as DC batteries, AC sine-wave generators, and Pulse sources, ensuring rigorous simulation accuracy.

\subsection{Junction and crossover detection}

Currently, the morphological wire-tracing algorithm assumes all intersections are shared electrical nodes. To support complex, overlapping schematics, "Junction Dots" and "Wire Crossovers" will be added as distinct detection classes. Furthermore, the OpenCV pipeline will utilize pixel-density heuristics to differentiate a genuine electrical connection from an insulated wire crossover.

\subsection{Expanding the component library}

To increase the utility of PySpice Studio as a robust analog design tool, the neural network's training dataset will be expanded. Rather than broadening into digital logic, focus will be placed on deeper analog integration—specifically training the model to recognize advanced MOSFET configurations and a wider array of specialized power sources. Synthetic data generation via programmatic drawing tools will be explored to rapidly scale this dataset without the bottleneck of manual human annotation.


% -------- 8. Conclusion --------
\section{Conclusion}

This phase of the research has successfully established a robust methodology for the digitization of hand-drawn circuit schematics into mathematically categorized device entities. By fine-tuning and deploying a YOLOv8-based object detection architecture, the proposed system effectively circumvents the inherent limitations of traditional computer vision algorithms, achieving highly resilient recognition performance across diverse and unconstrained sketching conventions. The critical milestone of precise spatial localization and multi-class component extraction has been achieved, establishing a solid technical foundation for subsequent topological wire-routing analysis. Full integration of the node-connectivity algorithm in future iterations will culminate in a comprehensive, end-to-end image-to-SPICE compiler, significantly streamlining the transition between conceptual hardware prototyping and rigorous software simulation.

% -------- References --------
\newpage
\addcontentsline{toc}{section}{References}
\begin{thebibliography}{99}

\makeatletter
\def\@biblabel#1{[#1]}
\makeatother

\section*{Research Papers}

\bibitem{redmon2016}
J. Redmon, S. Divvala, R. Girshick, and A. Farhadi, 
\href{https://www.cv-foundation.org/openaccess/content_cvpr_2016/papers/Redmon_You_Only_Look_CVPR_2016_paper.pdf}
{``You Only Look Once: Unified, Real-Time Object Detection,''} 
in \textit{Proc. CVPR}, 2016.

\bibitem{thoma2021}
F. Thoma, J. Bayer, Y. Li, and A. Dengel, 
\href{https://www.researchgate.net/publication/354357000_A_Public_Ground-Truth_Dataset_for_Handwritten_Circuit_Diagram_Images}
{``A Public Ground-Truth Dataset for Handwritten Circuit Diagram Images,''} 
in \textit{Proc. ICDAR}, 2021.

\bibitem{ahmed2025}
N. Ahmed et al., 
\href{https://www.researchgate.net/publication/388159586_Digitize-HCD_A_Dataset_for_Digitization_of_Handwritten_Circuit_Diagrams}
{``Digitize-HCD: A Dataset for Digitization of Handwritten Circuit Diagrams,''} 
\textit{Data in Brief}, vol. 59, 2025.

\bibitem{patare2016}
M. D. Patare and M. Joshi, 
\href{https://www.ijcaonline.org/archives/volume143/number3/patare-2016-ijca-910058.pdf}
{``Hand-drawn Digital Logic Circuit Component Recognition using SVM,''} 
\textit{International Journal of Computer Applications}, vol. 143, no. 3, pp. 24--28, 2016.

\section*{Datasets}

\bibitem{cghd}
J. Bayer, 
\href{https://www.kaggle.com/datasets/johannesbayer/cghd1152/data}
{``Circuit Graph Handwritten Dataset (CGHD1152),''} 
Kaggle Dataset, 2021.



\section*{Useful Websites and Documentation}

\bibitem{ultralyticsdocs}
Ultralytics, 
\href{https://docs.ultralytics.com/}
{``Ultralytics YOLO Documentation.''}

\bibitem{schemdrawdocs}
SchemDraw, 
\href{https://schemdraw.readthedocs.io/en/stable/}
{``SchemDraw Documentation.''}

\bibitem{netron}
Netron, 
\href{https://netron.app/}
{``Netron: Visualizer for Neural Network, Deep Learning, and Machine Learning Models.''}

\bibitem{ngspice}
NGSPICE, 
\href{https://ngspice.sourceforge.io/docs/ngspice-46-manual.pdf}
{``NGSPICE User's Manual Version 46.''}

\section*{GitHub Repositories}

\textbf{Our Project:}

\bibitem{projimage}
R. Kevat, Y. Badwaik, and D. Nayak, 
\href{https://github.com/Ridham19/image-to-ngspice}
{``Image-to-Ngspice: Handwritten Circuit to SPICE Pipeline.''} 

\bibitem{projspice}
R. Kevat, Y. Badwaik, and D. Nayak, 
\href{https://github.com/Ridham19/Spice-Playground}
{``Spice-Playground.''} 

\textbf{Useful Repositories:}

\bibitem{reposchemdraw}
C. Delker, 
\href{https://github.com/cdelker/schemdraw}
{``SchemDraw GitHub Repository.''}

\bibitem{repoultralytics}
Ultralytics, 
\href{https://github.com/ultralytics/ultralytics}
{``Ultralytics GitHub Repository.''}

\end{thebibliography}

\end{document}